% ======================================================================
% Sections 1–2 — Introduction and Background (expanded)
% Reuses material from submitted_version and previous_revision, retitled
% to the Stage 45.1 scientific story (graph-free baseline + consumption).
% ======================================================================
\section{Introduction}\label{sec:intro}

Accurate traffic forecasting is a core problem in intelligent transportation
systems: navigation, ride-hailing, route planning, and congestion management
all depend on reliable short-horizon speed or flow estimates
\citep{akhtar2021review}. The task is difficult because traffic state is
both \emph{spatial} --- downstream roads respond to upstream conditions, and
commuting couples distant corridors --- and \emph{temporal} --- patterns
evolve over minutes and hours. Graph-based deep models, especially Graph
Convolutional Networks (GCNs) \citep{kipf2017semi} and the temporal
extension T-GCN \citep{zhao2019t}, have become a standard tool: sensors are
nodes, an adjacency matrix $\mathbf{A}$ encodes spatial coupling, and a
recurrent or convolutional backbone maps a window of past speeds to future
speeds.

In almost all of these systems, $\mathbf{A}$ is \emph{given}. It is built
from road connectivity, sensor distance, or a heuristic kernel. That default
encodes a strong assumption: physical proximity approximates functional
traffic dependency. The assumption is often false. Distant roads can be
strongly coupled by commuter flows; adjacent sensors can be nearly
independent. A second, less discussed failure mode is \emph{oversmoothing}:
when $\mathbf{A}$ is dense, repeated graph convolution averages away
node-specific signal, so a model with a rich graph can be worse than the
same model with \emph{no} graph.

Graph structure learning (GSL) offers a data-driven alternative: estimate
$\mathbf{A}$ from the traffic series rather than from the map
\citep{Chen2022}. Continuous score-based DAG learners such as NOTEARS
\citep{zheng2018dags} and DAGMA \citep{Bello2024DAGMA} make this practical
by replacing combinatorial search with smooth optimization. Applying such
learners to traffic is attractive: the resulting sparse adjacency can be
read as a statistical dependency graph over sensors.

Yet two design questions remain open, and they organize this paper.

\begin{enumerate}
\item \textbf{What is the right reference graph?} If a dense physical graph
is harmful at the operating point of interest, then a learned graph that
merely \emph{improves on the physical baseline} may still be worse than
using the identity (graph-free) adjacency. Without a graph-free control,
GSL gains are not identified.
\item \textbf{How is the graph consumed over time?} A learned adjacency can
be applied as one static operator to the whole input window, or as
lag-specific operators aligned with input timesteps. Recurrent backbones
can use different operators at different steps; non-recurrent models
cannot. \emph{Structure} and \emph{consumption} are separate design
dimensions.
\end{enumerate}

A system bibliography of the field makes the first question concrete: a
survey of $784$ traffic-forecasting articles ($2000$--$2025$) shows that
graph-based methods have become dominant while still relying almost
universally on heuristic graph construction (Appendix~C). The literature
rarely reports a graph-free control, so it is hard to tell whether reported
GSL gains come from learned structure or from removing a harmful default.

\subsection*{This study}

We evaluate GCN and T-GCN backbones on two public benchmarks --- Los-loop
($207$ highway sensors, $5$\,min sampling) and SZ-Taxi ($156$ urban
sensors, $15$\,min sampling) --- under one protocol that crosses four
\emph{graph sources} with four \emph{consumption mechanisms}:
physical road-network graphs; identity (graph-free) controls; learned
\emph{contemporaneous} DAGMA graphs (GSL and its binary symmetrization
cGSL); and learned \emph{multi-lag} DAGMA graphs consumed by a fixed cyclic
assignment, a global lag mixture, a per-node per-timestep gate
(T-GCN-MultiGSL-Mix), or a static union inside GCN. We use four horizons,
five training seeds, and matched-sparsity and capacity controls.

\paragraph{Proposed method (summary).}
The technical core is the separation of \emph{estimating} lag-specific
dependency graphs from \emph{using} them (full development in
Sect.~\ref{sec:method}). We stack an explicit lag window
$Z=[x(t-L),\ldots,x(t)]$ and fit DAGMA once on that matrix, then extract
binary lag-block adjacencies. Those static graphs are supplied to T-GCN
variants that either cycle them across input steps (T-GCN-MultiGSL), mix
them with three global scalars (Weighted), or reweight them per node and
per timestep with a small gate (Mix). A GCN counterpart consumes the
element-wise union of the same graphs as a single static operator, providing
a controlled comparison of consumption on nearly identical edge sets.

\paragraph{Main findings.}
(i)~The physical graph is not a sound default: graph-free models beat it in
every dataset--horizon cell, on both backbones. (ii)~Single learned
contemporaneous graphs improve on the physical graph but never on the
graph-free control. (iii)~Multi-lag graphs help on Los-loop only when
consumed in a temporally aligned way; the identical graphs unioned into one
static adjacency fail. (iv)~Gated mixing is the strongest Los-loop
configuration (up to $14.5\%$ relative RMSE reduction versus the graph-free
baseline at PH1; $5/5$ paired seeds; up to $43.0\%$ versus the physical
T-GCN baseline across horizons). (v)~On SZ-Taxi, where the multi-lag fit
retains only two edges, all learned variants stay within seed variability of
graph-free; matched $30$-edge controls show sparsity alone does not explain
the Los-loop gain. Throughout, learned graphs are treated as \emph{static
statistical dependency estimates}, not causal maps.

\paragraph{Contributions.}
\begin{enumerate}
\item A graph-free-baseline reassessment of graph structure learning for
traffic forecasting on two public benchmarks.
\item A controlled dissociation indicating that graph consumption is a major
determinant of whether multi-lag structure is useful (with the backbone
confound stated explicitly).
\item A conditional positive result for gated per-timestep consumption of
lag-specific learned graphs (T-GCN-MultiGSL-Mix).
\item Matched-budget evidence that learned edge placement, not sparsity
alone, drives the remaining Los-loop gain.
\end{enumerate}

\paragraph{Organization.}
Section~\ref{sec:background} reviews background, including the DAGMA
estimator. Section~\ref{sec:method} formalizes the task, both graph-learning
constructions, and the consumption mechanisms. Section~\ref{sec:setup}
states the protocol. Section~\ref{sec:results} reports results.
Section~\ref{sec:discussion} discusses interpretation.
Section~\ref{sec:limitations} lists limitations.
Section~\ref{sec:conclusion} concludes.

\section{Background and Related Work}\label{sec:background}

\subsection{Traffic forecasting: a brief landscape}\label{sec:bg-landscape}

Classical approaches model each series or a small neighborhood with ARIMA,
SVR, Kalman filters, or shallow neural nets; they capture temporal
regularity but struggle with network-scale spatial coupling
\citep{Hamed1995Short,Smola2004A,Lippi2013Short,akhtar2021review}. Deep
spatio-temporal models (ST-ResNet, recurrent CNNs, LSTMs) improve temporal
capacity but remain Euclidean unless a graph is introduced explicitly
\citep{Zhang2016Deep,Cao2017Interactive}.

Graph-based methods dominate recent traffic forecasting. GCN-style
aggregation \citep{kipf2017semi} and T-GCN \citep{zhao2019t} are the
canonical lightweight backbones. Attention and transformer variants
(A3T-GCN, dynamic spatial--temporal attention, graph transformers) reweight
neighbors over time but still require an initial graph
\citep{A3T-GCN,Sun2023Attention}. In that sense they change \emph{how} a
given structure is used, not \emph{which} structure is the right one. This
paper deliberately isolates the structure/consumption question on GCN and
T-GCN so that architecture scale does not confound the comparison.

\subsection{Graph convolutional networks}\label{sec:bg-gcn}

A GCN layer maps node features $\mathbf{H}^{(l)}$ to
$\mathbf{H}^{(l+1)}$ by
\begin{equation}\label{eq:gcn-layer}
\mathbf{H}^{(l+1)}
=\sigma\!\left(
\tilde{\mathbf{D}}^{-\frac12}\tilde{\mathbf{A}}
\tilde{\mathbf{D}}^{-\frac12}
\mathbf{H}^{(l)}\mathbf{W}^{(l)}
\right),
\end{equation}
where $\tilde{\mathbf{A}}=\mathbf{A}+\mathbf{I}$ adds self-loops,
$\tilde{\mathbf{D}}$ is the degree matrix of $\tilde{\mathbf{A}}$, and
$\sigma$ is a nonlinearity \citep{kipf2017semi}. The normalized operator
$\hat{A}=\tilde{\mathbf{D}}^{-1/2}\tilde{\mathbf{A}}
\tilde{\mathbf{D}}^{-1/2}$ aggregates neighborhood information. In our
implementation every adjacency --- physical, identity, and learned --- is
passed through the same Laplacian-with-self-loops normalization
(Sect.~\ref{sec:task}), so the graph-free control is the identity operator
in that pipeline, not a missing-graph special case.

\subsection{Temporal graph convolutional networks}\label{sec:bg-tgcn}

T-GCN \citep{zhao2019t} embeds graph convolution inside a Gated Recurrent
Unit so that spatial aggregation and temporal gating interact at every
input step. Writing $\mathbf{Z}_t$ for the graph-convolved features at time
$t$, the GRU updates a hidden state $\mathbf{S}_t$ with update and reset
gates; forecasting is trained with an MSE-type loss plus an $\ell_2$
penalty on weights. T-GCN is the recurrent backbone for all multi-lag
consumption mechanisms in this paper. The non-recurrent GCN backbone applies
graph convolution to the flattened input window and produces the prediction
in a single step --- which is why it can only consume a \emph{static}
adjacency (Sect.~\ref{sec:consumption}).

\subsection{Graph structure learning and DAGMA}\label{sec:bg-dagma}

Graph structure learning estimates $\mathbf{A}$ (or a weighted $W$) from
data. In traffic, the motivation is that proximity is a poor proxy for
functional co-movement \citep{Chen2022}. Among GSL methods we use
\textbf{DAGMA} \citep{Bello2024DAGMA}, a continuous learner for directed
acyclic graphs that extends the NOTEARS program
\citep{zheng2018dags}.

\paragraph{From discrete DAGs to continuous optimization.}
Learning a DAG over $d$ variables by search is combinatorial. NOTEARS
replaces the discrete constraint with a smooth equality $h(W)=0$ on a real
matrix $W\in\mathbb{R}^{d\times d}$, so that gradient-based methods apply.
DAGMA strengthens this idea with a \emph{log-determinant} acyclicity
characterization based on $M$-matrices. For a sufficiently large scalar
$s>0$, define
\begin{equation}\label{eq:dagma-h}
h_{\mathrm{ldet}}(W)
= -\log\det\!\big(s\mathbf{I}-W\circ W\big) + d\log s,
\end{equation}
where $\circ$ is the Hadamard product. Then $h_{\mathrm{ldet}}(W)=0$ if and
only if the graph induced by $W$ is acyclic; positive values quantify
violation. Because $h_{\mathrm{ldet}}$ is smooth where defined, DAGMA can
optimize a data-fit score under (or penalized by) this constraint via an
augmented Lagrangian,
\begin{equation}\label{eq:dagma-al}
\mathcal{L}^{\rho}(W,\alpha)
= \mathrm{score}(W)
+ \tfrac{\rho}{2}\,\big|h_{\mathrm{ldet}}(W)\big|^{2}
+ \alpha\, h_{\mathrm{ldet}}(W),
\end{equation}
with dual variable $\alpha$ and penalty $\rho>0$.

\paragraph{Linear SEM score used here.}
We use the linear structural-equation variant (\texttt{DagmaLinear}) with
an $\ell_2$ loss and an $\ell_1$ sparsity penalty $\lambda_1$ on the
coefficients. After optimization, a binary adjacency is obtained by an
absolute-magnitude threshold on $|W_{ij}|$ (sign-symmetric, matching the
library semantics). In our code path the threshold regimes differ by
construction --- an internal fit threshold for the contemporaneous graphs
and a consumer threshold for the multi-lag graphs --- and are specified
exactly in Sects.~\ref{sec:gsl_contemp} and \ref{sec:gsl_multilag}.

\paragraph{Interpretation limits.}
Two cautions matter for the rest of the paper. First, acyclicity is an
\emph{optimization regularizer} that makes continuous learning tractable;
it is not evidence that traffic influence is acyclic, and we never claim
causal discovery. Second, the learned $W$ is a \emph{statistical}
dependency estimate under linearity and noise assumptions. Where we speak
of lag blocks, that reading is tied to how $Z$ is built
(Sect.~\ref{sec:gsl_multilag}), not to a validated physical propagation
model.

\subsection{Structure versus consumption}\label{sec:bg-consumption}

Much of the GSL-for-traffic literature treats a learned adjacency as a
drop-in replacement for a physical one: same backbone, new $\mathbf{A}$.
That framing silently assumes consumption is irrelevant. Recurrent models,
however, can apply different graph operators at different input steps, and
non-recurrent models cannot. This paper makes the distinction explicit:
three lag-specific DAGMA graphs are shared by T-GCN variants (fixed cyclic
assignment, global weights, or a per-node gate) and by a GCN that consumes
their static union. Empirically, that shared-graph design is what allows us
to ask whether utility attaches to the edge set or to the way it is used
(Sects.~\ref{sec:res-dissociation} and \ref{sec:disc-consumption}).

\subsection{Evaluation practice}\label{sec:bg-eval}

Reported GSL gains are usually relative to a physical or heuristic graph.
A graph-free identity control, matched-sparsity controls, and multi-seed
reporting are less common, but they are exactly what is needed if learned
structure is to be credited with gains that \emph{removing} a harmful graph
already provides. This study therefore treats NoSpatial as a first-class
baseline rather than an ablation curiosity.
